\documentclass[12pt,letterpaper]{article}
\usepackage[utf8]{inputenc}
\usepackage[spanish]{babel}
\usepackage{amsmath}
\usepackage{amsfonts}
\usepackage{amssymb}
\usepackage[margin=2.5cm]{geometry}

\begin{document}

\begin{center}
    \Large \textbf{Evaluación de Examen 2: Unidad II} \\
    \large Física para Ciencias de la Salud (005-1713) \\
    \vspace{0.5cm}
    \textbf{Estudiante:} Anamerc Acuña \quad \textbf{C.I.:} 33.279.962
    \vspace{0.3cm} \\
    \large \textbf{Calificación Final: 8/10 puntos}
\end{center}

\vspace{0.5cm}

\section*{Análisis de Resultados}

\subsection*{1. Cinemática (Aceleración media)}
\textbf{Procedimiento y resultado:} La estudiante plantea correctamente la ecuación general de la cinemática ($v_f = v_i + a \cdot t$) y ejecuta de manera adecuada la operación aritmética $\frac{0,6}{0,15} = 4$. Sin embargo, falla en el análisis dimensional al expresar el resultado final de la aceleración con unidades de velocidad ($4 \text{ m/s}$ en lugar de $4 \text{ m/s}^2$). \\
\textbf{Veredicto:} \textbf{Parcialmente correcto} (Penalización por unidades incorrectas) (1,5/2 pts).

\subsection*{2. Dinámica (Leyes de Newton)}
\textbf{Procedimiento y resultado:} Identifica el uso de la Segunda Ley de Newton ($\Sigma \vec{F} = m \cdot \vec{a}$), pero comete un \textbf{error crítico de álgebra} al despejar la incógnita. Escribe analíticamente $a = \frac{m}{F}$ en lugar de la forma correcta $a = \frac{F}{m}$. Esta falla procedimental la lleva a calcular $25 / 150 = 0,16$, además de acompañar el valor erróneo con unidades también incorrectas para una aceleración ($0,16 \text{ m/s}$). El resultado real del bloque era $6 \text{ m/s}^2$. \\
\textbf{Veredicto:} \textbf{Incorrecto} (Se otorgan puntos parciales mínimos por identificación de datos y ecuación base) (0,5/2 pts).

\subsection*{3. Trabajo Mecánico}
\textbf{Procedimiento y resultado:} Aplica de forma impecable la fórmula del trabajo al considerar la alineación de la fuerza y el desplazamiento ($\cos(0^\circ)$). El cálculo final producto de $(400 \text{ N}) \cdot (0,8 \text{ m})$ es perfectamente exacto: $320 \text{ J}$. \\
\textbf{Veredicto:} \textbf{Correcto} (2/2 pts).

\subsection*{4. Energía Potencial}
\textbf{Procedimiento y resultado:} Extrae los datos de masa y altura apropiadamente. Al desarrollar la fórmula de energía potencial gravitatoria ($E_p = m \cdot g \cdot h$), la estudiante asigna prematuramente la unidad Joules ($1,8 \text{ J}$) a un paso intermedio que aún no involucraba la gravedad; no obstante, el desarrollo global no pierde consistencia y el resultado final numérico de $17,64 \text{ J}$ es totalmente correcto. \\
\textbf{Veredicto:} \textbf{Correcto} (2/2 pts).

\subsection*{5. Potencia Mecánica}
\textbf{Procedimiento y resultado:} Plantea correctamente la relación entre trabajo y tiempo de ejecución. La sustitución es directa y la aritmética de $75 \text{ J} / 60 \text{ s}$ arroja la cifra exacta y verificada de $1,25 \text{ W}$. \\
\textbf{Veredicto:} \textbf{Correcto} (2/2 pts).

\vspace{0.5cm}
\section*{Comentarios Generales y Sugerencias}
El examen evidencia un buen dominio conceptual, sobre todo en los apartados de Energía, Trabajo y Potencia, acompañados de un formato limpio que incluye la sección de "Datos", "Fórmula Utilizada" y los desarrollos correspondientes.
\begin{itemize}
    \item Se aconseja practicar a fondo los **despejes algebraicos básicos**. En la física de ciencias de la salud, una variable mal despejada de una fracción (como ocurrió en la Pregunta 2) puede llevar a resultados analíticos completamente opuestos a la realidad.
    \item Es imperativo revisar siempre el **análisis dimensional**. Las magnitudes físicas deben corresponder exactamente a sus unidades en el S.I. (la aceleración nunca debe indicarse como $\text{m/s}$, puesto que alude a una velocidad).
\end{itemize}

\end{document}